\begin{codebox}
\codeline{\textcolor{cmtgray}{\#\ 概率本体推理示例}}
\codeline{\textcolor{cmtgray}{\#\ Probabilistic\ Ontology\ Reasoning\ Examples}}
\codeline{\textcolor{cmtgray}{\#}}
\codeline{\textcolor{cmtgray}{\#\ 经典描述逻辑只能表达"非真即假"的确定性知识}}
\codeline{\textcolor{cmtgray}{\#\ 概率本体在公理上附加概率度量，用于处理工程中的不确定性}}
\codeblank{}
\codeline{\textcolor{cmtgray}{\#\ ==============================}}
\codeline{\textcolor{cmtgray}{\#\ 1.\ 概率类包含公理}}
\codeline{\textcolor{cmtgray}{\#\ ==============================}}
\codeline{\textcolor{cmtgray}{\#\ 在\ TBox\ 公理上标注条件概率（P-SROIQ\ /\ PR-OWL\ 风格）}}
\codeblank{}
\codeline{\textcolor{cmtgray}{\#\ 高功率设备需要维护的概率为85\%：}}
\codeline{P(NeedsMaintenance\ |\ HighPowerEquipment)\ =\ 0.85}
\codeline{\textcolor{cmtgray}{\#\ 含义：100台高功率设备中，约85台需要定期维护}}
\codeblank{}
\codeline{\textcolor{cmtgray}{\#\ 长时间运行的设备出现精度漂移的概率为60\%：}}
\codeline{P(PrecisionDrift\ |\ LongRunning)\ =\ 0.60}
\codeblank{}
\codeline{\textcolor{cmtgray}{\#\ ==============================}}
\codeline{\textcolor{cmtgray}{\#\ 2.\ 贝叶斯网络条件概率表\ (CPT)}}
\codeline{\textcolor{cmtgray}{\#\ ==============================}}
\codeline{\textcolor{cmtgray}{\#\ 设备故障诊断贝叶斯网络结构：}}
\codeline{\textcolor{cmtgray}{\#}}
\codeline{\textcolor{cmtgray}{\#\ \ \ HighTemp（高温）\ \ \ \ \ LongRunning（长时间运行）}}
\codeline{\textcolor{cmtgray}{\#\ \ \ \ \ \ \ \ \ \ \textbackslash{}\ \ \ \ \ \ \ \ \ \ \ \ \ \ \ \ /}}
\codeline{\textcolor{cmtgray}{\#\ \ \ \ \ \ \ \ \ \ \ v\ \ \ \ \ \ \ \ \ \ \ \ \ \ v}}
\codeline{\textcolor{cmtgray}{\#\ \ \ \ \ \ \ \ \ \ \ \ \ Fault（故障）}}
\codeline{\textcolor{cmtgray}{\#}}
\codeline{\textcolor{cmtgray}{\#\ 条件概率表：}}
\codeline{P(Fault\ |\ HighTemp\ \(\land\)\ LongRunning)\ \ \ =\ 0.90}
\codeline{P(Fault\ |\ HighTemp\ \(\land\)\ \(\lnot\)LongRunning)\ \ =\ 0.30}
\codeline{P(Fault\ |\ \(\lnot\)HighTemp\ \(\land\)\ LongRunning)\ \ =\ 0.10}
\codeline{P(Fault\ |\ \(\lnot\)HighTemp\ \(\land\)\ \(\lnot\)LongRunning)\ =\ 0.01}
\codeblank{}
\codeline{\textcolor{cmtgray}{\#\ 先验概率：}}
\codeline{P(HighTemp)\ =\ 0.20\ \ \ \ \ \ \#\ 车间环境中设备高温的先验概率}
\codeline{P(LongRunning)\ =\ 0.35\ \ \ \#\ 设备长时间运行的先验概率}
\codeblank{}
\codeline{\textcolor{cmtgray}{\#\ ==============================}}
\codeline{\textcolor{cmtgray}{\#\ 3.\ 概率推理计算示例}}
\codeline{\textcolor{cmtgray}{\#\ ==============================}}
\codeline{\textcolor{cmtgray}{\#\ 预测方向：观测到\ Lathe\_003\ 高温且长时间运行，故障概率多大？}}
\codeline{\textcolor{cmtgray}{\#\ 直接查条件概率表：P(Fault)\ =\ 0.90\ \(\rightarrow\)\ 触发预防性维护工单}}
\codeblank{}
\codeline{\textcolor{cmtgray}{\#\ 诊断方向：已知某设备发生故障，它处于高温状态的概率多大？}}
\codeline{\textcolor{cmtgray}{\#\ 第一步，用全概率公式计算故障的边缘概率：}}
\codeline{P(Fault)\ =\ 0.90\(\times\)0.20\(\times\)0.35\ +\ 0.30\(\times\)0.20\(\times\)0.65}
\codeline{\hspace*{9\ttcw}+\ 0.10\(\times\)0.80\(\times\)0.35\ +\ 0.01\(\times\)0.80\(\times\)0.65}
\codeline{\hspace*{9\ttcw}=\ 0.063\ +\ 0.039\ +\ 0.028\ +\ 0.0052}
\codeline{\hspace*{9\ttcw}=\ 0.1352}
\codeblank{}
\codeline{\textcolor{cmtgray}{\#\ 第二步，计算高温条件下的故障概率：}}
\codeline{P(Fault\ |\ HighTemp)\ =\ 0.90\(\times\)0.35\ +\ 0.30\(\times\)0.65\ =\ 0.51}
\codeblank{}
\codeline{\textcolor{cmtgray}{\#\ 第三步，用贝叶斯公式反推：}}
\codeline{P(HighTemp\ |\ Fault)\ =\ P(Fault\ |\ HighTemp)\ \(\times\)\ P(HighTemp)\ /\ P(Fault)}
\codeline{\hspace*{20\ttcw}=\ (0.51\ \(\times\)\ 0.20)\ /\ 0.1352}
\codeline{\hspace*{20\ttcw}\(\approx\)\ 0.754}
\codeblank{}
\codeline{\textcolor{cmtgray}{\#\ 结论：故障设备中约75\%与高温相关\ \(\rightarrow\)\ 维修时优先排查冷却系统}}
\codeblank{}
\codeline{\textcolor{cmtgray}{\#\ ==============================}}
\codeline{\textcolor{cmtgray}{\#\ 4.\ 马尔可夫逻辑网络\ (MLN)\ 权重规则}}
\codeline{\textcolor{cmtgray}{\#\ ==============================}}
\codeline{\textcolor{cmtgray}{\#\ MLN\ 为一阶逻辑规则附加权重：权重越大，约束越"硬"}}
\codeline{\textcolor{cmtgray}{\#\ 权重语义：满足该规则的可能世界，其概率按\ exp(权重)\ 倍提升}}
\codeblank{}
\codeline{2.5\ \ HighPowerEquipment(x)\ \(\rightarrow\)\ NeedsCooling(x)\ \ \ \ \ \ \#\ 强约束}
\codeline{1.2\ \ CNCMachine(x)\ \(\land\)\ Aged(x)\ \(\rightarrow\)\ PrecisionDrift(x)\ \ \#\ 中等约束}
\codeline{0.4\ \ Idle(x)\ \(\rightarrow\)\ Available(x)\ \ \ \ \ \ \ \ \ \ \ \ \ \ \ \ \ \ \ \ \ \ \ \#\ 弱约束（默认倾向）}
\codeblank{}
\codeline{\textcolor{cmtgray}{\#\ 对比：硬约束（等价于无穷大权重）直接用\ OWL\ 公理表达}}
\codeline{Running\ \(\sqcap\)\ Maintenance\ \(\sqsubseteq\)\ \(\bot\)\ \ \ \#\ 运行中的设备不可同时维护（确定性）}
\codeblank{}
\codeline{\textcolor{cmtgray}{\#\ ==============================}}
\codeline{\textcolor{cmtgray}{\#\ 5.\ 工程应用与工具}}
\codeline{\textcolor{cmtgray}{\#\ ==============================}}
\codeline{\textcolor{cmtgray}{\#\ 适用场景：}}
\codeline{\textcolor{cmtgray}{\#\ \ \ 设备故障诊断与预测性维护（PHM）}}
\codeline{\textcolor{cmtgray}{\#\ \ \ 质量风险评估（工艺参数\ \(\rightarrow\)\ 缺陷概率）}}
\codeline{\textcolor{cmtgray}{\#\ \ \ 多源传感器数据融合（不确定信息合并）}}
\codeline{\textcolor{cmtgray}{\#}}
\codeline{\textcolor{cmtgray}{\#\ 代表工具\ /\ 语言：}}
\codeline{\textcolor{cmtgray}{\#\ \ \ PR-OWL\ \ \ ——\ 基于\ MEBN\ 的概率\ OWL\ 扩展}}
\codeline{\textcolor{cmtgray}{\#\ \ \ BayesOWL\ ——\ 将\ OWL\ 本体转换为贝叶斯网络}}
\codeline{\textcolor{cmtgray}{\#\ \ \ Pronto\ \ \ ——\ Pellet\ 推理器的概率扩展（P-SROIQ）}}
\codeline{\textcolor{cmtgray}{\#\ \ \ ProbLog\ /\ Alchemy\ ——\ 概率逻辑编程\ /\ MLN\ 实现}}
\codeline{\textcolor{cmtgray}{\#}}
\codeline{\textcolor{cmtgray}{\#\ 实践建议：保持"确定性骨架\ +\ 概率标注"分离——}}
\codeline{\textcolor{cmtgray}{\#\ 类层次与硬约束用\ OWL\ DL\ 表达，不确定性知识在概率层叠加，}}
\codeline{\textcolor{cmtgray}{\#\ 避免将全部知识概率化导致推理复杂度失控}}
\end{codebox}
